\chapter{Reducción de endomorfismos}\label{ch:b2:reduction}

Para entender un endomorfismo hay que hallar las direcciones que se
limita a dilatar. Este capítulo construye la maquinaria
---valores propios, \omterm{def:b2:reduction:charpoly}{polinomios característico} y
mínimo, lema de descomposición en núcleos--- y recoge sus frutos:
criterios de diagonalización y de trigonalización, Cayley--Hamilton,
la \omterm{thm:b2:reduction:dunford}{descomposición de Dunford} y el cálculo de potencias y
exponenciales del que se alimentará el~\cref{ch:b2:diffeq}. En todo el
capítulo, $E$ es un $K$-espacio vectorial de dimensión finita ($K = \R$
o $\C$), $u \in \mathcal{L}(E)$ y $n = \dim E$.

\section{Valores y vectores propios}

\begin{definition}\label{def:b2:reduction:eigen}
$\lambda \in K$ es un \emph{valor propio}\index{valor propio} de $u$
cuando $u(x) = \lambda x$ para algún $x \neq 0$ (un \emph{vector
propio}); el \emph{subespacio propio} es $E_\lambda(u) = \ker(u -
\lambda\,\mathrm{id})$. El conjunto de los valores propios es el
\emph{espectro}\index{espectro} $\operatorname{Sp}(u)$. Un subespacio
$F$ es \emph{estable} cuando $u(F) \subseteq F$; los subespacios
propios son estables, y los subespacios estables permiten definir
endomorfismos inducidos $u|_F$.
\end{definition}

\begin{theorem}[Independencia de los subespacios propios]\label{thm:b2:reduction:independence}
Los \omterm{def:b2:reduction:eigen}{vectores propios} asociados a \omterm{def:b2:reduction:eigen}{valores propios} distintos dos a
dos forman una familia libre; equivalentemente, la suma de los
\omterm{def:b2:reduction:eigen}{subespacios propios} $E_{\lambda_1} + \dots + E_{\lambda_r}$ (con
$\lambda_i$ distintos) es directa. En particular, $u$ tiene a lo sumo
$n$ \omterm{def:b2:reduction:eigen}{valores propios}.
\end{theorem}

\begin{proof}
Por inducción sobre $r$. Supongamos $x_1 + \dots + x_r = 0$ con $x_i
\in E_{\lambda_i}$, conocido el enunciado para $r - 1$. Apliquemos $u$
y restemos $\lambda_r$ veces la relación:
\[
\sum_{i=1}^{r-1} (\lambda_i - \lambda_r)\, x_i = 0 ,
\]
así que por inducción cada $(\lambda_i - \lambda_r)x_i = 0$, es decir,
$x_i = 0$ para $i < r$, y entonces $x_r = 0$. Una suma directa de
espacios no nulos dentro de un espacio de dimensión $n$ tiene a lo
sumo $n$ sumandos.
\end{proof}

\begin{omfigure}
\begin{tikzpicture}[scale=0.9, >={Stealth}]
  \draw[gray!40] (-2.6,-2.2) grid (2.6,2.2);
  \draw[->] (-2.6,0) -- (2.7,0) node[below, font=\small] {$x$};
  \draw[->] (0,-2.2) -- (0,2.3) node[left, font=\small] {$y$};
  \draw[omDef, very thick, ->] (0,0) -- (0.7,0.7)
    node[above, font=\small] {$v_1$};
  \draw[omDef, thick, ->, dashed] (0,0) -- (2.1,2.1)
    node[above right, font=\small] {$Av_1 = 3v_1$};
  \draw[omThm, very thick, ->] (0,0) -- (0.9,-0.9)
    node[below, font=\small] {$v_2$};
  \draw[omThm, thick, ->, dashed] (0,0) -- (0.9,-0.9);
  \draw[omProp, very thick, ->] (0,0) -- (1,0)
    node[below right, font=\small] {$e_1$};
  \draw[omProp, thick, ->, dashed] (0,0) -- (2,1)
    node[right, font=\small] {$Ae_1$};
\end{tikzpicture}

{\small La matriz $A = \begin{psmallmatrix}2 & 1\\ 1 &
2\end{psmallmatrix}$ actuando sobre el plano: el vector genérico
$e_1$ sale despedido de su recta, pero las direcciones propias $v_1 =
(1,1)$ y $v_2 = (1,-1)$ se limitan a dilatarse, por $3$ y por $1$
(así que $Av_2 = v_2$: la imagen discontinua coincide con $v_2$).
Diagonalizar es cambiar a la base $(v_1, v_2)$, donde $A$ pasa a ser
$\operatorname{diag}(3, 1)$.}
\end{omfigure}

\begin{definition}[Polinomio característico]\label{def:b2:reduction:charpoly}
$\chi_u(X) = \det(X\,\mathrm{id} - u)$\index{polinomio
característico}, calculado en cualquier base como $\det(XI_n - A)$;
es un polinomio mónico de grado $n$, invariante por semejanza
(\cref{thm:b2:linalg:detrules}). Sus raíces en $K$ son exactamente
los \omterm{def:b2:reduction:eigen}{valores propios} ($\lambda$ es \omterm{def:b2:reduction:eigen}{valor propio} $\iff u -
\lambda\,\mathrm{id}$ no es inyectiva $\iff \chi_u(\lambda) = 0$), y
\[
\chi_u(X) = X^n - (\operatorname{tr} u)\, X^{n-1} + \dots +
(-1)^n \det u .
\]
La \emph{multiplicidad algebraica}\index{multiplicidad algebraica}
$m_\lambda$ de un \omterm{def:b2:reduction:eigen}{valor propio} es su multiplicidad como raíz de
$\chi_u$; la \emph{multiplicidad geométrica}\index{multiplicidad
geométrica} es $\dim E_\lambda$, y se cumple $1 \leq \dim E_\lambda
\leq m_\lambda$.
\end{definition}

\begin{proof}[Demostración de los hechos enunciados]
Sobre los coeficientes: desarróllese $\det(XI - A)$ por la fórmula de
las permutaciones; la permutación identidad aporta $\prod_i (X -
a_{ii}) = X^n - (\sum a_{ii})X^{n-1} + \dots$, y cualquier otra
permutación deja fijas a lo sumo $n - 2$ posiciones diagonales, con
lo que aporta grado $\leq n - 2$: los dos coeficientes superiores son
los indicados; y $X = 0$ da el término constante $\det(-A) =
(-1)^n\det A$.

Geométrica $\leq$ algebraica: sea $d = \dim E_\lambda$ y complétese
una base de $E_\lambda$ hasta una base de $E$; la matriz de $u$ es
triangular superior por bloques con bloque superior izquierdo
$\lambda I_d$, luego $\chi_u(X) = (X -
\lambda)^d\, \chi_{\text{(bloque inferior)}}(X)$: la multiplicidad de
$\lambda$ es al menos $d$.
\end{proof}

\begin{example}[Mismo $\chi$, geometría distinta]\label{ex:b2:reduction:samechi}
Las matrices
\[
\begin{pmatrix}2 & 0\\ 0 & 2\end{pmatrix}
\qquad\text{y}\qquad
\begin{pmatrix}2 & 1\\ 0 & 2\end{pmatrix}
\]
comparten el \omterm{def:b2:reduction:charpoly}{polinomio característico} $(X - 2)^2$, la traza, el
\omterm{def:b2:linalg:det}{determinante} y el \omterm{def:b2:reduction:eigen}{espectro}, y sin embargo no son semejantes: la
primera tiene $E_2$ de dimensión $2$ (\omterm{def:b2:reduction:charpoly}{multiplicidad geométrica}
$2$) y la segunda de dimensión $1$. El \omterm{def:b2:reduction:charpoly}{polinomio característico}
solo ve las \omterm{def:b2:reduction:charpoly}{multiplicidades algebraicas}; las dimensiones de los
\omterm{def:b2:reduction:eigen}{subespacios propios} son el invariante más fino, y el
\omterm{def:b2:reduction:polyu}{polinomio mínimo} es quien arbitra ($X - 2$ frente a $(X - 2)^2$).
Moraleja para toda discusión sobre \omterm{def:b2:reduction:diag}{diagonalizabilidad}: $\chi$
preselecciona a los candidatos, pero quienes votan son los núcleos.
\end{example}

\begin{definition}[Diagonalizable, trigonalizable]\label{def:b2:reduction:diag}
$u$ es \emph{diagonalizable}\index{diagonalizable} cuando $E$ tiene
una base de \omterm{def:b2:reduction:eigen}{vectores propios} (en términos matriciales: semejante a
una matriz diagonal); es \emph{trigonalizable}\index{trigonalizable}
cuando su matriz en alguna base es triangular superior.
\end{definition}

\begin{theorem}[Criterios de diagonalizabilidad]\label{thm:b2:reduction:diagcrit}
Las afirmaciones siguientes son equivalentes:
\begin{enumerate}
  \item $u$ es \omterm{def:b2:reduction:diag}{diagonalizable};
  \item $E = \bigoplus_{\lambda \in \operatorname{Sp} u} E_\lambda$;
  \item $\chi_u$ se escinde sobre $K$ y $\dim E_\lambda = m_\lambda$
        para todo \omterm{def:b2:reduction:eigen}{valor propio};
  \item (suficiente, no necesaria) $\chi_u$ tiene $n$ raíces
        distintas en $K$.
\end{enumerate}
\end{theorem}

\begin{proof}
(1 $\iff$ 2): una base de \omterm{def:b2:reduction:eigen}{vectores propios} se reparte en bases de
los $E_\lambda$, y recíprocamente, concatenando bases de los sumandos
directos se obtiene una base de $E$
(el~\cref{thm:b2:reduction:independence} hace la suma directa; la
igualdad de dimensiones hace que lo llene todo).

(2 $\iff$ 3): en la base diagonal, $\chi_u = \prod (X -
\lambda)^{\dim E_\lambda}$ se escinde con multiplicidades que
coinciden. Recíprocamente, supongamos que $\chi_u$ se escinde con
$\dim E_\lambda = m_\lambda$ en todos los casos; entonces la suma
directa de los \omterm{def:b2:reduction:eigen}{subespacios propios} (directa por el~\cref{thm:b2:reduction:independence}) tiene dimensión
\[
\sum_{\lambda}\dim E_\lambda = \sum_{\lambda} m_\lambda =
\deg\chi_u = n ,
\]
donde la igualdad central se debe a que el grado de un polinomio
escindido es la suma de las multiplicidades de sus raíces: la
suma es todo $E$. Obsérvese dónde ha trabajado cada hipótesis: la
escisión ha llenado el grado y la igualdad de multiplicidades ha
llenado las dimensiones.

(4 $\Rightarrow$ 1): $n$ \omterm{def:b2:reduction:eigen}{valores propios} distintos dan $n$
\omterm{def:b2:reduction:eigen}{vectores propios} independientes
(\cref{thm:b2:reduction:independence}): una base.
\end{proof}

\begin{method}[Cómo decidir la diagonalizabilidad]\label{met:b2:reduction:decide}
En la práctica conviene comprobar en este orden, pues cada paso
puede zanjar el asunto. (1) ¿Se presenta por sí solo un polinomio
\omterm{def:b2:linalg:annihilator}{anulador} escindido con raíces simples ($u^2 = \mathrm{id}$, $u^2 =
u$, $u^k = \mathrm{id}$)? Si es así: \omterm{def:b2:reduction:diag}{diagonalizable}, sin ningún
cálculo (\cref{cor:b2:reduction:minpolycrit} más abajo). (2)
Calcúlese $\chi_u$; si tiene $n$ raíces distintas en $K$:
\omterm{def:b2:reduction:diag}{diagonalizable} (\cref{thm:b2:reduction:diagcrit}~(4)). (3) En caso
contrario, y solo para cada raíz múltiple $\lambda$, compárese
$\dim\ker(u - \lambda\,\mathrm{id})$ con la multiplicidad
$m_\lambda$: cualquier déficit mata la \omterm{def:b2:reduction:diag}{diagonalizabilidad}, y la
igualdad en todos los casos la demuestra. Nunca hay que calcular
los \omterm{def:b2:reduction:eigen}{subespacios propios} de las raíces simples (su dimensión está
forzada a ser $1$), ni trigonalizar solo para decidir.
\end{method}

\begin{example}[La diagonalización puesta a trabajar]\label{ex:b2:reduction:diagwork}
$A = I + J = \begin{psmallmatrix}2 & 1 & 1\\ 1 & 2 & 1\\ 1 & 1 &
2\end{psmallmatrix}$, con $J$ la matriz de unos: de
$\operatorname{Sp}(J) = \{3, 0\}$
(\cref{ex:b2:linalg:onesmatrix}) se sigue $\operatorname{Sp}(A) =
\{4, 1\}$, con \omterm{def:b2:reduction:eigen}{subespacios propios} $\R(1,1,1)$ y el plano $\{x + y
+ z = 0\}$: dimensiones $1 + 2 = 3$, luego \omterm{def:b2:reduction:diag}{diagonalizable}
(\cref{thm:b2:reduction:diagcrit}~(2)). Potencias sin ninguna matriz
de cambio de base: con $\Pi = J/3$ el proyector sobre $\R(1,1,1)$,
\[
A = 4\,\Pi + 1\cdot(I - \Pi)
\quad\Longrightarrow\quad
A^k = 4^k\,\Pi + (I - \Pi)
= \frac{4^k - 1}{3}\,J + I .
\]
(Comprobación con $k = 1$: $\frac{4-1}3 J + I = A$.) La moraleja:
cuando los \omterm{def:b2:reduction:eigen}{subespacios propios} son visibles, los \emph{proyectores}
espectrales calculan potencias más deprisa de lo que jamás lo hará
$PDP^{-1}$, y además la fórmula muestra la dinámica: $A^k$ crece como
$4^k$ a lo largo de $(1,1,1)$ y se queda quieto en el plano
ortogonal.
\end{example}

\begin{theorem}[Trigonalización]\label{thm:b2:reduction:trigonalization}
$u$ es \omterm{def:b2:reduction:diag}{trigonalizable} sobre $K$ si y solo si $\chi_u$ se escinde
sobre $K$. En particular, todo endomorfismo de un $\C$-espacio
vectorial es \omterm{def:b2:reduction:diag}{trigonalizable}.\index{trigonalización}
\end{theorem}

\begin{proof}
($\Rightarrow$) El \omterm{def:b2:reduction:charpoly}{polinomio característico} de una matriz triangular
es $\prod(X - t_{ii})$: escindido.

($\Leftarrow$) Por inducción sobre $n$. Como $\chi_u$ se escinde,
tiene una raíz $\lambda$: tomemos un \omterm{def:b2:reduction:eigen}{vector propio} $e_1$. En una
base que empiece por $e_1$, la matriz es $\begin{pmatrix} \lambda &
\ast\\ 0 & B\end{pmatrix}$, y $\chi_u = (X - \lambda)\chi_B$, así que
$\chi_B$ también se escinde. Por la hipótesis de inducción aplicada a
la matriz $B$ de tamaño $(n-1) \times (n-1)$, existe una $Q$
invertible con $Q^{-1}BQ$ triangular superior; conjugar la matriz
entera por $\begin{pmatrix}1 & 0\\ 0 & Q\end{pmatrix}$ la
triangulariza.
\end{proof}

\begin{example}[Trigonalizando a mano]\label{ex:b2:reduction:trigwork}
$B = \begin{pmatrix}3 & -1\\ 1 & 1\end{pmatrix}$: $\chi_B = X^2 -
4X + 4 = (X - 2)^2$, y $\ker(B - 2I) = \ker\begin{psmallmatrix}1 &
-1\\ 1 & -1\end{psmallmatrix}$ es la recta \omterm{def:b2:structures:generated}{generada} por $e_1' = (1,
1)$: un solo \omterm{def:b2:reduction:eigen}{valor propio} y un \omterm{def:b2:reduction:eigen}{subespacio propio} de dimensión
uno; no es \omterm{def:b2:reduction:diag}{diagonalizable}, pero sí \omterm{def:b2:reduction:diag}{trigonalizable}
(\cref{thm:b2:reduction:trigonalization}). Completemos la base con
$e_2' = (1, 0)$ y calculemos:
\[
u(e_1') = (2, 2) = 2e_1',
\qquad
u(e_2') = (3, 1) = 1\cdot e_1' + 2\, e_2' ,
\]
de modo que en la base $(e_1', e_2')$ la matriz es $T =
\begin{psmallmatrix}2 & 1\\ 0 & 2\end{psmallmatrix}$. La moraleja:
la diagonal de $T$ estaba forzada (ambas entradas han de ser el
\omterm{def:b2:reduction:eigen}{valor propio} doble $2$); solo la entrada de la esquina dependía de
la elección de $e_2'$, y reescalar $e_2'$ permite darle cualquier
valor no nulo. Ese ``$1$'' que se resiste es la sombra de la parte
nilpotente que Dunford aislará.
\end{example}

\section{Polinomios de un endomorfismo}

\begin{definition}\label{def:b2:reduction:polyu}
Para $P = \sum a_k X^k \in K[X]$, pongamos $P(u) = \sum a_k u^k \in
\mathcal{L}(E)$. La aplicación $P \mapsto P(u)$ es un morfismo de
álgebras $K[X] \to \mathcal{L}(E)$
(\cref{def:b2:structures:algebra}); su núcleo $\{P : P(u) = 0\}$ es
un \omterm{def:b2:structures:ideal}{ideal} de $K[X]$, no nulo (la familia $(\mathrm{id}, u, \dots,
u^{n^2})$ está ligada en $\mathcal{L}(E)$, de dimensión $n^2$), y por
tanto está \omterm{def:b2:structures:generated}{generado} por un único polinomio mónico $\mu_u$: el
\emph{polinomio mínimo}\index{polinomio mínimo}
(\cref{thm:b2:structures:principal}).
\end{definition}

\begin{proposition}\label{prop:b2:reduction:minpoly}
\begin{enumerate}
  \item $P(u) = 0 \iff \mu_u \mid P$; los \omterm{def:b2:reduction:eigen}{valores propios} de $u$ son
        raíces de todo polinomio \omterm{def:b2:linalg:annihilator}{anulador}, y las raíces de $\mu_u$
        son \emph{exactamente} los \omterm{def:b2:reduction:eigen}{valores propios}.
  \item Si $F$ es estable, entonces $\mu_{u|_F} \mid \mu_u$.
\end{enumerate}
\end{proposition}

\begin{proof}
(1) La divisibilidad es la definición de generador. Si $u(x) =
\lambda x$ con $x \neq 0$, entonces $0 = P(u)(x) = P(\lambda)x$,
luego $P(\lambda) = 0$: los \omterm{def:b2:reduction:eigen}{valores propios} son raíces de los
\omterm{def:b2:linalg:annihilator}{anuladores} y, en particular, de $\mu_u$. Recíprocamente, si
$\lambda$ es raíz, $\mu_u = (X - \lambda)Q$ con $Q(u) \neq 0$ (el
grado de $\mu_u$ es mínimo): tómese $y$ con $Q(u)(y) \neq 0$;
entonces $(u - \lambda)(Q(u)(y)) = \mu_u(u)(y) = 0$ exhibe el
\omterm{def:b2:reduction:eigen}{vector propio} $Q(u)(y)$.

(2) $\mu_u(u|_F) = \mu_u(u)|_F = 0$, y se aplica (1) a $u|_F$.
\end{proof}

\begin{example}[Polinomios mínimos hallados a mano]\label{ex:b2:reduction:minpolyhand}
El \omterm{def:b2:reduction:polyu}{polinomio mínimo} se calcula probando grados sucesivos. Para
la matriz de unos $J \in \mathcal{M}_3(\R)$: $J \neq \lambda I$
(queda descartado el grado $1$) y $J^2 = 3J$, luego
\[
\mu_J = X^2 - 3X = X(X - 3) :
\]
grado $2$, escindido, con raíces simples; $J$ es \omterm{def:b2:reduction:diag}{diagonalizable}
con \omterm{def:b2:reduction:eigen}{espectro} $\{0, 3\}$ (\cref{cor:b2:reduction:minpolycrit} más
abajo), lo que confirma el~\cref{ex:b2:linalg:onesmatrix} sin
calcular ni un solo \omterm{def:b2:linalg:det}{determinante}. Para la matriz de intercambio $A$
del~\cref{ex:b2:reduction:projectorswork}: de $A \neq \pm I$ y $A^2
= I$ resulta $\mu_A = X^2 - 1$. En ambos casos el patrón es el
mismo: adivínese a partir de la estructura una identidad de grado
bajo (el rango uno obliga a $J^2 = (\operatorname{tr}J)\,J$; una
involución obliga a $A^2 = I$) y compruébese después que ningún
divisor propio anula. Los \omterm{def:b2:reduction:polyu}{polinomios mínimos} suelen
\emph{encontrarse}, no calcularse a partir de $\chi$.
\end{example}

\begin{theorem}[Lema de descomposición en núcleos]\label{thm:b2:reduction:kernels}
Si $P = P_1 P_2 \cdots P_r$ con los $P_i$ primos entre sí dos a dos,
entonces\index{lema de descomposición en núcleos}
\[
\ker P(u) = \ker P_1(u) \oplus \dots \oplus \ker P_r(u),
\]
y las proyecciones sobre los sumandos son polinomios en $u$.
\end{theorem}

\begin{proof}
Basta tratar el caso $r = 2$ e inducir. Por Bézout en $K[X]$
(\cref{thm:b2:structures:principal}), $U P_1 + V P_2 = 1$, luego para
todo $x$,
\[
x = \underbrace{U(u)P_1(u)(x)}_{=:\,x_2}
+ \underbrace{V(u)P_2(u)(x)}_{=:\,x_1}.
\]
Si $x \in \ker P(u)$, entonces $P_2(u)(x_2) = U(u)\,P(u)(x) = 0$ (los
polinomios en $u$ conmutan), luego $x_2 \in \ker P_2(u)$, y
simétricamente $x_1 \in \ker P_1(u)$: la suma llena $\ker P(u)$; y
ambos sumandos están dentro de $\ker P(u)$ (pues $P_i \mid P$).
Carácter directo: si $x \in \ker P_1(u) \cap \ker P_2(u)$, entonces
$x = U(u)P_1(u)x + V(u)P_2(u)x = 0$. Las fórmulas para $x_1, x_2$
exhiben las proyecciones como $V(u)P_2(u)$ y $U(u)P_1(u)$.
\end{proof}

\begin{example}[El lema de los núcleos con proyectores explícitos]\label{ex:b2:reduction:projectorswork}
Sea $A = \begin{psmallmatrix}0 & 1 & 0\\ 1 & 0 & 0\\ 0 & 0 &
1\end{psmallmatrix}$ (intercambia las dos primeras coordenadas).
Entonces $A^2 = I$: el polinomio $X^2 - 1 = (X - 1)(X + 1)$ anula
$A$, sus factores son primos entre sí y Bézout es explícito:
\[
\frac{1}{2}(X + 1) - \frac12(X - 1) = 1 .
\]
Siguiendo la demostración del~\cref{thm:b2:reduction:kernels}, las
proyecciones sobre $\ker(A - I)$ y $\ker(A + I)$ son los polinomios
en $A$
\[
\pi_+ = \frac{A + I}{2} = \frac12\begin{pmatrix}
1 & 1 & 0\\ 1 & 1 & 0\\ 0 & 0 & 2\end{pmatrix},
\qquad
\pi_- = \frac{I - A}{2} = \frac12\begin{pmatrix}
1 & -1 & 0\\ -1 & 1 & 0\\ 0 & 0 & 0\end{pmatrix}.
\]
Comprobación: $\pi_+ + \pi_- = I$, $\pi_+\pi_- = 0$, $\pi_\pm^2 =
\pi_\pm$, y las imágenes son el plano $\{x = y\}$ (vectores
simétricos, \omterm{def:b2:reduction:eigen}{valor propio} $1$) y la recta $\R(1, -1, 0)$
(antisimétricos, \omterm{def:b2:reduction:eigen}{valor propio} $-1$). El lema de los núcleos no es
un enunciado de existencia: los coeficientes de Bézout \emph{son}
las fórmulas de los proyectores.
\end{example}

\begin{example}[Los proyectores calculan también la exponencial]\label{ex:b2:reduction:expprojectors}
La misma matriz de intercambio, un dividendo más allá. Como $A =
\pi_+ - \pi_-$ con proyectores ortogonales en sentido algebraico
($\pi_+\pi_- = 0$), toda potencia cumple $A^k = \pi_+ +
(-1)^k\pi_-$, y la serie exponencial se reagrupa por proyectores:
\[
\eu^{tA} = \sum_k \frac{t^k}{k!}\bigl(\pi_+ +
(-1)^k\pi_-\bigr)
= \eu^{t}\,\pi_+ + \eu^{-t}\,\pi_- =
\begin{pmatrix}
\cosh t & \sinh t & 0\\
\sinh t & \cosh t & 0\\
0 & 0 & \eu^{t}
\end{pmatrix}.
\]
(Comprobación en $t = 0$: la identidad; derivada en $0$: $A$.) La
descomposición espectral convierte una serie de matrices en dos
series escalares, que es exactamente el mecanismo que el~\cref{ch:b2:diffeq} aplicará a todo sistema \omterm{def:b2:reduction:diag}{diagonalizable}, y la
razón de que las funciones hiperbólicas gobiernen los acoplamientos
simétricos.
\end{example}

\begin{corollary}[Diagonalizabilidad mediante el polinomio mínimo]\label{cor:b2:reduction:minpolycrit}
$u$ es \omterm{def:b2:reduction:diag}{diagonalizable} $\iff$ $\mu_u$ se escinde sobre $K$ con
raíces \emph{simples} $\iff$ algún polinomio \omterm{def:b2:linalg:annihilator}{anulador} de $u$ se
escinde con raíces simples.
\end{corollary}

\begin{proof}
Si $P(u) = 0$ con $P = \prod_{i}(X - \lambda_i)$ (los $\lambda_i$
distintos), el lema da $E = \ker P(u) = \bigoplus_i \ker(u -
\lambda_i)$: una suma directa de \omterm{def:b2:reduction:eigen}{subespacios propios}, luego $u$ es
\omterm{def:b2:reduction:diag}{diagonalizable} (\cref{thm:b2:reduction:diagcrit}). Recíprocamente,
un $u$ \omterm{def:b2:reduction:diag}{diagonalizable} es anulado por $\prod_{\lambda \in
\operatorname{Sp}u}(X - \lambda)$ (que anula cada \omterm{def:b2:reduction:eigen}{subespacio
propio}), polinomio escindido con raíces simples; y $\mu_u$ lo
divide teniendo las mismas raíces
(\cref{prop:b2:reduction:minpoly}): $\mu_u$ es exactamente ese
producto.
\end{proof}

\begin{example}\label{ex:b2:reduction:projections}
Las proyecciones cumplen $p^2 = p$: las anula $X(X-1)$, escindido
con raíces simples, luego son \omterm{def:b2:reduction:diag}{diagonalizables} con \omterm{def:b2:reduction:eigen}{espectro}
$\subseteq \{0, 1\}$, y $E = \ker p \oplus \ker(p - \mathrm{id})$: el
análisis geométrico del primer año, vuelto a demostrar en una línea.
Las simetrías ($s^2 = \mathrm{id}$, \omterm{def:b2:linalg:annihilator}{anulador} $X^2 - 1$) son
\omterm{def:b2:reduction:diag}{diagonalizables} cuando $\operatorname{char} K \neq 2$, con
\omterm{def:b2:reduction:eigen}{espectro} $\subseteq \{\pm 1\}$. Un endomorfismo con $u^3 = u^2$ y
$u^2 \neq u$ está anulado por $X^2(X - 1)$ y \emph{no} es
necesariamente \omterm{def:b2:reduction:diag}{diagonalizable}: el criterio lo detecta (hay que
examinar la raíz doble $0$; es \omterm{def:b2:reduction:diag}{diagonalizable} si y solo si además
$\ker u^2 = \ker u$).
\end{example}

\begin{example}[El cuerpo decide: una rotación en $\R^3$]\label{ex:b2:reduction:rotation3d}
Sea $R$ el cuarto de vuelta alrededor del eje $z$:
\[
R = \begin{pmatrix}
0 & -1 & 0\\
1 & 0 & 0\\
0 & 0 & 1
\end{pmatrix},
\qquad
\chi_R = (X - 1)(X^2 + 1).
\]
Sobre $\R$: el único \omterm{def:b2:reduction:eigen}{valor propio} es $1$, con \omterm{def:b2:reduction:eigen}{subespacio propio} el
eje $\R e_3$ ---una sola recta de vectores fijos y ninguna reducción
más---: $R$ no es \omterm{def:b2:reduction:diag}{diagonalizable} ni \omterm{def:b2:reduction:diag}{trigonalizable} en
$\mathcal{M}_3(\R)$ (pues $\chi_R$ no se escinde). Sobre $\C$: tres
\omterm{def:b2:reduction:eigen}{valores propios} distintos $1, \iu, -\iu$, luego $R$ es
\omterm{def:b2:reduction:diag}{diagonalizable}, con \omterm{def:b2:reduction:eigen}{vectores propios} $e_3$ y $e_1 \mp \iu e_2$.
La geometría se oía ya en el álgebra: las rotaciones del plano no
tienen direcciones invariantes reales, y los \omterm{def:b2:reduction:eigen}{valores propios}
complejos $\pm\iu$, de módulo $1$, guardan el ángulo
($\pm\frac\pi2$) que la matriz real solo puede expresar mezclando
coordenadas.
\end{example}

\begin{example}[Mínimo frente a característico]\label{ex:b2:reduction:minvschar}
Para $D = \operatorname{diag}(2, 2, 3)$: $\chi_D = (X - 2)^2(X -
3)$, pero $\mu_D = (X - 2)(X - 3)$, ya que $(D - 2I)(D - 3I) = 0$
(compruébese sobre la base canónica) mientras que ninguno de los dos
factores anula $D$ por separado. Para el bloque de desplazamiento $N
= \begin{psmallmatrix}0 & 1\\ 0 & 0\end{psmallmatrix} \oplus (3)$,
es decir, $N' = \begin{psmallmatrix}0 & 1 & 0\\ 0 & 0 & 0\\ 0 & 0 &
3\end{psmallmatrix}$: se tiene $\chi_{N'} = X^2(X - 3)$ \emph{y}
$\mu_{N'} = X^2(X - 3)$; la raíz doble hace verdadera falta porque
$N'$ no es \omterm{def:b2:reduction:diag}{diagonalizable} del lado del núcleo ($N'e_2 = e_1 \neq
0$). Regla práctica: $\mu$ y $\chi$ comparten sus raíces
(\cref{prop:b2:reduction:minpoly}); la multiplicidad en $\mu$ mide
el tamaño del mayor bloque nilpotente, y la de $\chi$ la dimensión
total del subespacio característico.
\end{example}

\begin{theorem}[Cayley--Hamilton]\label{thm:b2:reduction:cayleyhamilton}
$\chi_u(u) = 0$; en consecuencia, $\mu_u \mid \chi_u$ y $\deg \mu_u
\leq n$.\index{teorema de Cayley--Hamilton}
\end{theorem}

\begin{proof}
Fijemos $x \neq 0$ y sea $d$ máximo tal que $(x, u(x), \dots,
u^{d-1}(x))$ sea libre; escribamos
\[
u^d(x) = -a_0 x - a_1 u(x) - \dots - a_{d-1}u^{d-1}(x),
\]
y pongamos $P_x = X^d + a_{d-1}X^{d-1} + \dots + a_0$, de modo que
$P_x(u)(x) = 0$. Completemos la familia libre hasta una base de $E$:
en ella, $u$ tiene forma por bloques $\begin{pmatrix} C & \ast\\ 0 &
D\end{pmatrix}$, donde $C$ es la matriz compañera de $P_x$, cuyo
\omterm{def:b2:reduction:charpoly}{polinomio característico} es $P_x$ (desarróllese $\det(XI - C)$ por
la primera columna, por inducción sobre $d$). Por tanto $\chi_u = P_x
\cdot \chi_D$, y
\[
\chi_u(u)(x) = \chi_D(u)\bigl(P_x(u)(x)\bigr) = 0 .
\]
El argumento vale para todo $x$: $\chi_u(u) = 0$.
\end{proof}

\begin{example}[Cayley--Hamilton en acción]\label{ex:b2:reduction:chwork}
$A = \begin{pmatrix} 1 & 2\\ 3 & 4\end{pmatrix}$: $\chi_A = X^2 - 5X
- 2$, luego $A^2 = 5A + 2I$. Toda potencia de $A$ se reduce a una
combinación de $I$ y $A$:
\[
A^4 = (5A + 2I)^2 = 25A^2 + 20A + 4I = 145A + 54I =
\begin{pmatrix} 199 & 290\\ 435 & 634\end{pmatrix},
\]
y la inversa sale gratis: de $A(A - 5I) = 2I$ se obtiene
\[
A^{-1} = \tfrac12(A - 5I)
= \begin{pmatrix} -2 & 1\\ 3/2 & -1/2\end{pmatrix}.
\]
La moraleja: Cayley--Hamilton comprime toda el álgebra $K[A]$ en
$\operatorname{Vect}(I, A, \dots, A^{n-1})$ ---se tiene $\dim K[A] =
\deg\mu_A \leq n$, por grandes que sean las potencias que
necesites---.
\end{example}

\begin{remark}[Errores frecuentes]\label{rem:b2:reduction:pitfalls}
(i) Los \omterm{def:b2:reduction:eigen}{valores propios} no se suman: $\operatorname{Sp}(A + B)$ no
es $\operatorname{Sp}A + \operatorname{Sp}B$, y una suma de matrices
\omterm{def:b2:reduction:diag}{diagonalizables} no tiene por qué ser \omterm{def:b2:reduction:diag}{diagonalizable};
$\begin{psmallmatrix}1 & 1\\ 0 & 0\end{psmallmatrix} +
\begin{psmallmatrix}0 & 0\\ 0 & 1\end{psmallmatrix} =
\begin{psmallmatrix}1 & 1\\ 0 & 1\end{psmallmatrix}$ es suma de dos
matrices \omterm{def:b2:reduction:diag}{diagonalizables} (cada una con \omterm{def:b2:reduction:eigen}{valores propios}
distintos) y no es \omterm{def:b2:reduction:diag}{diagonalizable}; solo se portan bien las
familias que \emph{conmutan} (\cref{exo:b2:reduction:9}). (ii)
``$\chi_u$ se escinde'' es una hipótesis sobre el \emph{cuerpo}: una
rotación del plano tiene $\chi = X^2 - 2\cos\theta\,X + 1$,
escindido sobre $\C$ pero no sobre $\R$; es \omterm{def:b2:reduction:diag}{diagonalizable} en
$\mathcal{M}_2(\C)$ y ni siquiera \omterm{def:b2:reduction:diag}{trigonalizable} en
$\mathcal{M}_2(\R)$. (iii) La desigualdad va de geométrica $\leq$
algebraica, nunca al revés; comprobar solo que $\dim E_\lambda \geq
1$ no demuestra nada sobre la \omterm{def:b2:reduction:diag}{diagonalizabilidad}. (iv) $\mu_u$ no
es $\chi_u$: la igualdad se da exactamente cuando cada \omterm{def:b2:reduction:eigen}{valor propio}
tiene una única cadena de bloques (por ejemplo, las matrices
compañeras del problema de fin de semana de este capítulo); usar
$\chi$ donde hace falta $\mu$ infla todos los cálculos de potencias.
(v) El $d$ y el $\nu$ de Dunford son polinomios en $u$: una
descomposición $u = d' + \nu'$ con las propiedades adecuadas pero
con $d'\nu' \neq \nu'd'$ \emph{no} es la de Dunford y nunca es
única.
\end{remark}

\begin{remark}[Dónde se usa este capítulo]
La reducción es el caballo de batalla del resto del libro: las
potencias y exponenciales de matrices mueven los sistemas
diferenciales lineales del~\cref{ch:b2:diffeq}; el teorema espectral
del~\cref{ch:b2:quadratic} es la diagonalización hecha ortogonal;
las funciones generatrices (\cref{ch:b2:genfun}) vuelven a deducir
analíticamente las asintóticas de recurrencias del problema de fin
de semana de este capítulo. En el volumen del tercer año el mismo
programa se ejecuta en dimensión infinita: la teoría espectral de
los operadores compactos autoadjuntos, donde sucesiones de
\omterm{def:b2:reduction:eigen}{valores propios} sustituyen a los \omterm{def:b2:reduction:eigen}{espectros} finitos, y la teoría de
Perron--Frobenius de las matrices positivas, que explica \emph{por
qué} los \omterm{def:b2:reduction:eigen}{valores propios} dominantes de los problemas de recuento
son positivos y simples.
\end{remark}

\section{Nilpotentes y descomposición de Dunford}

\begin{proposition}[Endomorfismos nilpotentes]\label{prop:b2:reduction:nilpotent}
Para $u$ con $\chi_u$ escindido, las afirmaciones siguientes son
equivalentes: $u^n = 0$; $u^k = 0$ para algún $k$;
$\operatorname{Sp}(u) = \{0\}$; $\chi_u = X^n$; $u$ es
\omterm{def:b2:reduction:diag}{trigonalizable} con diagonal nula. Un \omterm{prop:b2:reduction:nilpotent}{endomorfismo nilpotente}
tiene $\mu_u = X^{\text{(índice de nilpotencia)}}$, con índice $\leq
n$.\index{endomorfismo nilpotente}
\end{proposition}

\begin{proof}
Si $u^k = 0$, todo \omterm{def:b2:reduction:eigen}{valor propio} es raíz de $X^k$: el \omterm{def:b2:reduction:eigen}{espectro} es
$\{0\}$ (no vacío cuando $\chi$ se escinde, y sobre $\C$ siempre).
Entonces $\chi_u = X^n$ (todas las raíces nulas) y Cayley--Hamilton
da $u^n = 0$; la trigonalización
(\cref{thm:b2:reduction:trigonalization}) pone ceros en la diagonal
(la diagonal lleva los \omterm{def:b2:reduction:eigen}{valores propios}). Recíprocamente, sea $A$
estrictamente triangular superior: $a_{ij} = 0$ para $j \leq i$.
Veamos por inducción que
\[
(A^k)_{ij} = 0 \qquad \text{siempre que } j \leq i + k - 1,
\]
es decir, que cada potencia empuja la región de ceros una diagonal
más arriba. Para $k = 1$ es la hipótesis. Para el paso inductivo,
\[
(A^{k+1})_{ij} = \sum_{\ell} (A^k)_{i\ell}\,a_{\ell j} ,
\]
y cada término se anula: o bien $\ell \leq i + k - 1$ (el primer
factor es $0$ por inducción), o bien $\ell \geq i + k$, en cuyo caso
$j \leq i + k \leq \ell$ mata el segundo factor. Para $k = n$, la
condición $j \leq i + n - 1$ se cumple para todos los $i, j \leq n$:
$A^n = 0$. El \omterm{def:b2:reduction:polyu}{polinomio mínimo} divide a $X^n$ y la anulación
define el índice.
\end{proof}

\begin{theorem}[Descomposición de Dunford]\label{thm:b2:reduction:dunford}
Supongamos que $\chi_u$ se escinde sobre $K$ (automático si $K =
\C$). Entonces existe un \emph{único} par $(d, \nu)$ con
\[
u = d + \nu, \qquad d \text{ diagonalizable}, \quad \nu
\text{ nilpotente}, \quad d\nu = \nu d ,
\]
y además $d$ y $\nu$ son polinomios en
$u$.\index{descomposición de Dunford}
\end{theorem}

\begin{proof}
\emph{Existencia.} Escribamos $\chi_u = \prod_{i=1}^{r} (X -
\lambda_i)^{m_i}$ (los $\lambda_i$ distintos) y pongamos $N_i =
\ker(u - \lambda_i)^{m_i}$, los \emph{subespacios característicos}.
Por Cayley--Hamilton y el lema de los núcleos
(\cref{thm:b2:reduction:kernels}),
\[
E = N_1 \oplus \dots \oplus N_r ,
\]
con proyecciones $\pi_i$ polinómicas en $u$; cada $N_i$ es estable
(los polinomios en $u$ conmutan con $u$). Definamos $d = \sum_i
\lambda_i \pi_i$: es un polinomio en $u$ y es \omterm{def:b2:reduction:diag}{diagonalizable} (actúa
como $\lambda_i$ sobre $N_i$, de modo que $E$ se descompone en sus
\omterm{def:b2:reduction:eigen}{subespacios propios}). Entonces $\nu = u - d$ es un polinomio en $u$
(y por tanto conmuta con $d$) y actúa sobre cada $N_i$ como $u -
\lambda_i$, con $(u - \lambda_i)^{m_i} = 0$ allí: $\nu^{\max m_i} =
0$ en cada sumando, luego $\nu$ es nilpotente.

\emph{Unicidad.} Sea $u = d' + \nu'$ otro par en esas condiciones.
Como $d'$ y $\nu'$ conmutan entre sí, conmutan con $u = d' + \nu'$ y
por tanto con todo polinomio en $u$; en particular, con $d$ y con
$\nu$. Entonces $d - d'$ es \omterm{def:b2:reduction:diag}{diagonalizable} (dos aplicaciones
\omterm{def:b2:reduction:diag}{diagonalizables} que conmutan son simultáneamente
\omterm{def:b2:reduction:diag}{diagonalizables}: \cref{exo:b2:reduction:9}) e igual a $\nu' - \nu$,
que es nilpotente: si $\nu^k = 0$ y $\nu'^{k'} = 0$, la conmutación
autoriza el desarrollo del binomio
\[
(\nu' - \nu)^{k + k' - 1}
= \sum_{j=0}^{k+k'-1}\binom{k + k' - 1}{j}
\,\nu'^{\,j}\,(-\nu)^{k + k' - 1 - j} ,
\]
en el que todos los términos mueren: o bien $j \geq k'$ (primer
factor nulo), o bien $k + k' - 1 - j \geq k$ (segundo factor
nulo), y siempre se da una de las dos. Un nilpotente
\omterm{def:b2:reduction:diag}{diagonalizable} es nulo (su \omterm{def:b2:reduction:eigen}{espectro} es $\{0\}$ y es diagonal en
alguna base): luego $d = d'$ y $\nu = \nu'$.
\end{proof}

\begin{example}[Potencias y exponenciales]\label{ex:b2:reduction:powers}
$A = \begin{pmatrix} 3 & 1\\ -1 & 1\end{pmatrix}$: $\chi_A = X^2 -
4X + 4 = (X-2)^2$, con un único \omterm{def:b2:reduction:eigen}{valor propio} $2$ y \omterm{def:b2:reduction:eigen}{subespacio
propio} de dimensión $1$: no es \omterm{def:b2:reduction:diag}{diagonalizable}. Dunford: $D = 2I$,
$N = A - 2I = \begin{pmatrix} 1 & 1\\ -1 & -1\end{pmatrix}$, con
$N^2 = 0$. Entonces
\[
A^k = (2I + N)^k = 2^k I + k\,2^{k-1} N ,
\qquad
\eu^{tA} = \eu^{2t}(I + tN),
\]
por el teorema del binomio para elementos que conmutan y por la
serie exponencial (\cref{ch:b2:diffeq}) separada sobre sumandos que
conmutan. La reducción convierte la dinámica matricial en dinámica
escalar.
\end{example}

\begin{remark}[Perspectivas dentro de este volumen]\label{rem:b2:reduction:perspectives}
La reducción es un nudo de comunicaciones; conviene vigilar cuatro
ramales. En el~\cref{ch:b2:nvs}, las normas adaptadas convierten
``todos los \omterm{def:b2:reduction:eigen}{valores propios} de módulo $< 1$'' en ``alguna norma de
operador $< 1$'', con lo que los \omterm{def:b2:reduction:eigen}{espectros} gobiernan la convergencia
de potencias y series. En el~\cref{ch:b2:diffeq}, la receta del~\cref{ex:b2:reduction:powers} pasa a ser la solución general de $X'
= AX$: Dunford separa $\eu^{tA}$ en bloques de tipo polinomio por
exponencial, y la estabilidad se lee en las partes reales de los
\omterm{def:b2:reduction:eigen}{valores propios}. En el~\cref{ch:b2:quadratic}, un producto escalar
impone lo que el álgebra lineal por sí sola no puede: las matrices
simétricas resultan \emph{ortogonalmente} \omterm{def:b2:reduction:diag}{diagonalizables}, sin parte
nilpotente alguna. Y en el~\cref{ch:b2:genfun}, las asintóticas por
\omterm{def:b2:reduction:eigen}{valor propio} dominante del problema de fin de semana de este
capítulo reaparecen analíticamente, como la singularidad más pequeña
de una función generatriz: dos lenguajes para una misma tasa de
crecimiento.
\end{remark}

\section{Ejercicios}

\begin{exercise}[$\star$]\label{exo:b2:reduction:1}
Diagonaliza (\omterm{def:b2:reduction:eigen}{valores propios}, bases de los \omterm{def:b2:reduction:eigen}{subespacios propios},
matriz invertible $P$):
\[
A = \begin{pmatrix} 1 & 2\\ 2 & 1 \end{pmatrix},
\qquad
B = \begin{pmatrix} 0 & 1 & 1\\ 1 & 0 & 1\\ 1 & 1 & 0
\end{pmatrix}.
\]
\end{exercise}

\begin{exercise}[$\star$]\label{exo:b2:reduction:2}
Prueba que $C = \begin{pmatrix} 1 & 1\\ 0 & 1\end{pmatrix}$ no es
\omterm{def:b2:reduction:diag}{diagonalizable} de dos maneras: mediante los \omterm{def:b2:reduction:eigen}{subespacios propios}
y mediante el \omterm{def:b2:reduction:polyu}{polinomio mínimo}.
\end{exercise}

\begin{exercise}[$\star$]\label{exo:b2:reduction:3}
Sea $u$ tal que $u^2 - 5u + 6\,\mathrm{id} = 0$. Demuestra que $u$ es
\omterm{def:b2:reduction:diag}{diagonalizable}, determina los \omterm{def:b2:reduction:eigen}{espectros} posibles y calcula $u^k$
como combinación de $\mathrm{id}$ y $u$.
\end{exercise}

\begin{exercise}[$\star\star$]\label{exo:b2:reduction:4}
Sea $u$ \omterm{def:b2:reduction:diag}{diagonalizable} y $F$ un subespacio estable. Demuestra que
$u|_F$ es \omterm{def:b2:reduction:diag}{diagonalizable} \emph{(restringe un polinomio \omterm{def:b2:linalg:annihilator}{anulador}
escindido con raíces simples)}.
\end{exercise}

\begin{exercise}[$\star\star$]\label{exo:b2:reduction:5}
(Fibonacci) Sea $A = \begin{pmatrix} 1 & 1\\ 1 & 0\end{pmatrix}$.
Diagonaliza $A$ sobre $\R$ y deduce la fórmula de Binet para la
sucesión de Fibonacci ($F_0 = 0$, $F_1 = 1$, $F_{n+1} = F_n +
F_{n-1}$):
\[
F_n = \frac{\varphi^n - \psi^n}{\sqrt 5},
\qquad \varphi = \frac{1 + \sqrt5}{2},\ \psi = \frac{1 -
\sqrt5}{2}.
\]
\end{exercise}

\begin{exercise}[$\star\star$]\label{exo:b2:reduction:6}
Sea $u \in \mathcal{L}(E)$ con $u^2$ \omterm{def:b2:reduction:diag}{diagonalizable} y $u$
invertible ($K = \C$). Demuestra que $u$ es \omterm{def:b2:reduction:diag}{diagonalizable}. Da un
contraejemplo cuando $u$ no es invertible.
\end{exercise}

\begin{exercise}[$\star\star$]\label{exo:b2:reduction:7}
Calcula la \omterm{thm:b2:reduction:dunford}{descomposición de Dunford}, $A^k$ y $\eu^{tA}$ para
\[
A = \begin{pmatrix} 2 & 1 & 0\\ 0 & 2 & 1\\ 0 & 0 & 2
\end{pmatrix}.
\]
\end{exercise}

\begin{exercise}[$\star\star$]\label{exo:b2:reduction:8}
Sea $A \in \mathcal{M}_n(\C)$ con $A^k = I$ para algún $k \geq 1$.
Demuestra que $A$ es \omterm{def:b2:reduction:diag}{diagonalizable} y que sus \omterm{def:b2:reduction:eigen}{valores propios} son
raíces $k$-ésimas de la unidad. Deduce que una matriz compleja
invertible de orden finito semejante a una matriz triangular con
diagonal de unos es la identidad.
\end{exercise}

\begin{exercise}[$\star\star\star$]\label{exo:b2:reduction:9}
(Diagonalización simultánea) Sean $u, v$ \omterm{def:b2:reduction:diag}{diagonalizables} y que
conmutan. Demuestra que son \emph{simultáneamente}
\omterm{def:b2:reduction:diag}{diagonalizables}: alguna base diagonaliza ambos. \emph{(Cada
\omterm{def:b2:reduction:eigen}{subespacio propio} de $u$ es estable por $v$; diagonaliza allí las
restricciones de $v$ usando el~\cref{exo:b2:reduction:4}.)}
\end{exercise}

\begin{exercise}[$\star\star\star$]\label{exo:b2:reduction:10}
Sea $u \in \mathcal{L}(\C^n)$. Demuestra que $u$ es
\omterm{def:b2:reduction:diag}{diagonalizable} si y solo si todo subespacio estable por $u$ admite
un suplementario estable por $u$. \emph{(Para $\Leftarrow$: aplica la
propiedad a $F = \sum_\lambda E_\lambda(u)$, la suma de todos los
\omterm{def:b2:reduction:eigen}{subespacios propios}; si un suplementario estable $G$ fuera no
nulo, trigonalizar $u|_G$ produciría un \omterm{def:b2:reduction:eigen}{vector propio} de $u$ dentro
de $G$, en contradicción con $G \cap F = \{0\}$.)}
\end{exercise}

\begin{exercise}[$\star\star\star$]\label{exo:b2:reduction:11}
(Radio espectral en versión ligera de Gelfand, un anticipo $2\times2$
del análisis que viene) Sea $A \in \mathcal{M}_2(\C)$ con ambos
\omterm{def:b2:reduction:eigen}{valores propios} de módulo $< 1$. Demuestra que $A^k \to 0$ entrada
a entrada cuando $k \to \infty$. \emph{(Trigonaliza: $A =
P(T)P^{-1}$ con $T$ triangular superior; calcula $T^k$
explícitamente ---distinguiendo \omterm{def:b2:reduction:eigen}{valores propios} iguales y
distintos--- y acota.)}
\end{exercise}

\begin{exercise}[$\star\star$]\label{exo:b2:reduction:12}
Sea $u \in \mathcal{L}(\C^n)$ con $\operatorname{rk} u = 1$ ($n \geq
2$). Prueba que $\chi_u = X^{n-1}(X - \operatorname{tr} u)$ y que $u$
es \omterm{def:b2:reduction:diag}{diagonalizable} si y solo si $\operatorname{tr} u \neq 0$.
\emph{(Recuerda del~\cref{exo:b2:linalg:5} que $u^2 =
(\operatorname{tr} u)\,u$.)}
\end{exercise}

\section{Problema: recurrencias lineales y matrices compañeras}

Una recurrencia lineal $u_{n+k} = a_{k-1}u_{n+k-1} + \dots + a_0 u_n$
es una potencia de matriz disfrazada, y la reducción la convierte en
fórmulas cerradas, tasas de crecimiento y estimaciones del error. Este
problema de fin de semana desarrolla el diccionario ---matrices
compañeras de un lado, operador de desplazamiento sobre el espacio de
las sucesiones del otro---, demuestra el \emph{teorema fundamental de
las recurrencias lineales} (la solución general es $\sum_i
Q_i(n)\lambda_i^n$ sobre las raíces del \omterm{def:b2:reduction:charpoly}{polinomio característico}) y
gasta los dividendos en la aproximación diofántica de $\sqrt2$, en el
recuento de caminos y palabras, y en un anillo de sucesiones acopladas
que solo la diagonalización simultánea consigue desenredar.

\begin{problem}[{Problema de fin de semana --- el teorema fundamental
de las recurrencias lineales}]
\label{pb:b2:reduction:1}
Fijemos $k \geq 1$, escalares $a_0, \dots, a_{k-1} \in \C$ con $a_0
\neq 0$, el polinomio mónico $P = X^k - a_{k-1}X^{k-1} - \dots - a_1
X - a_0$ y la recurrencia
\[
(\mathcal R)\colon\quad u_{n+k} = a_{k-1}u_{n+k-1} + \dots +
a_1 u_{n+1} + a_0 u_n \qquad (n \geq 0).
\]
La \emph{matriz compañera} de $P$ es
\[
C =
\begin{pmatrix}
0 & 1 & & \\
 & \ddots & \ddots & \\
 & & 0 & 1\\
a_0 & a_1 & \cdots & a_{k-1}
\end{pmatrix}
\in \mathcal{M}_k(\C).
\]

\medskip
\noindent\textbf{Parte I --- El diccionario de la matriz
compañera.}
\begin{enumerate}
  \item Prueba que una sucesión $(u_n)$ satisface $(\mathcal R)$ si
        y solo si los vectores $v_n = (u_n, u_{n+1}, \dots,
        u_{n+k-1})^{\mathsf T}$ satisfacen $v_{n+1} = Cv_n$, de
        donde $v_n = C^n v_0$.
  \item Demuestra que $\chi_C = P$ (desarrolla $\det(XI - C)$ por
        la primera columna e induce sobre $k$) y después que
        también $\mu_C = P$ \emph{(pasa a $C^{\mathsf T}$, para la
        que $e_1$ es \omterm{def:b2:structures:generated}{cíclico}, y observa que una matriz y su
        \omterm{def:b2:linalg:transpose}{traspuesta} tienen el mismo \omterm{def:b2:reduction:polyu}{polinomio mínimo})}.
  \item Prueba que, para cada raíz $\lambda$ de $P$, el vector $(1,
        \lambda, \dots, \lambda^{k-1})^{\mathsf T}$ genera el
        \omterm{def:b2:reduction:eigen}{subespacio propio} de $C$ asociado a $\lambda$; deduce que
        \emph{todos} los \omterm{def:b2:reduction:eigen}{subespacios propios} de $C$ tienen
        dimensión $1$ y que $C$ es \omterm{def:b2:reduction:diag}{diagonalizable} si y solo si $P$
        tiene $k$ raíces distintas.
  \item Supongamos que $P$ tiene raíces distintas $\lambda_1, \dots,
        \lambda_k$. Prueba que las sucesiones geométricas
        $(\lambda_i^n)_n$ forman una base del espacio de soluciones
        de $(\mathcal R)$, de modo que toda solución es $u_n =
        \sum_i c_i\lambda_i^n$ para constantes $c_i$ únicas.
  \item Resuelve por completo: $u_{n+2} = u_{n+1} + 6u_n$, $u_0 =
        1$, $u_1 = 8$.
\end{enumerate}

\medskip
\noindent\textbf{Parte II --- El operador de desplazamiento y el
teorema fundamental.} Sea $\mathcal{S}$ el $\C$-espacio vectorial de
todas las sucesiones complejas y $S \in \mathcal{L}(\mathcal{S})$ el
desplazamiento, $S\bigl((u_n)_n\bigr) = (u_{n+1})_n$.
\begin{enumerate}[resume]
  \item Prueba que el conjunto de soluciones de $(\mathcal R)$ es
        $\ker P(S)$ y que tiene dimensión exactamente $k$
        \emph{(envía cada solución a sus valores iniciales)}.
  \item Explica por qué el lema de descomposición en núcleos
        (\cref{thm:b2:reduction:kernels}) se aplica a $S$ sobre el
        espacio $\mathcal{S}$, de dimensión infinita, sin ningún
        cambio, y escribe la descomposición resultante de $\ker
        P(S)$ para $P = \prod_{i=1}^{r}(X - \lambda_i)^{m_i}$ (los
        $\lambda_i$ distintos, todos no nulos porque $a_0 \neq 0$).
  \item Para $\lambda \neq 0$ y $m \geq 1$, prueba que
        \[
        \ker\,(S - \lambda\,\mathrm{id})^m
        = \bigl\{\,\bigl(Q(n)\,\lambda^n\bigr)_n : Q \in
        \C_{m-1}[X]\,\bigr\},
        \]
        de dimensión $m$. \emph{(Calcula $(S -
        \lambda)\bigl(Q(n)\lambda^n\bigr) =
        \lambda^{n+1}(\Delta Q)(n)$ con $\Delta Q = Q(X + 1) -
        Q(X)$, y usa que $\Delta$ baja el grado; para la dimensión,
        acótala por $m$ mediante los valores iniciales.)}
  \item (El teorema fundamental de las recurrencias lineales)
        Concluye: si $P = \prod_{i=1}^{r}(X - \lambda_i)^{m_i}$ con
        los $\lambda_i$ distintos y no nulos, las soluciones de
        $(\mathcal R)$ son exactamente las sucesiones
        \[
        u_n = \sum_{i=1}^{r} Q_i(n)\,\lambda_i^n,
        \qquad Q_i \in \C_{m_i - 1}[X],
        \]
        con polinomios $Q_i$ determinados de manera única.
  \item Resuelve por completo: $u_{n+2} = 4u_{n+1} - 4u_n$, $u_0 =
        1$, $u_1 = 0$, y comprueba la respuesta sobre $u_2$.
\end{enumerate}

\medskip
\noindent\textbf{Parte III --- Raíces dominantes y dividendos
diofánticos.}
\begin{enumerate}[resume]
  \item Supongamos las raíces simples con $\abs{\lambda_1} >
        \abs{\lambda_i}$ para $i \geq 2$, y $u_n = \sum_i c_i
        \lambda_i^n$ con $c_1 \neq 0$. Prueba que $u_n \sim
        c_1\lambda_1^n$ y que $u_{n+1}/u_n \to \lambda_1$.
  \item (Pell) Definamos $a_{n+1} = a_n + 2b_n$, $b_{n+1} = a_n +
        b_n$, $a_0 = b_0 = 1$. Prueba que $q(a, b) = a^2 - 2b^2$
        cumple $q(a_{n+1}, b_{n+1}) = -q(a_n, b_n)$, de donde
        $a_n^2 - 2b_n^2 = (-1)^{n+1}$; relaciónalo con el
        \omterm{def:b2:linalg:det}{determinante} de $M = \begin{psmallmatrix}1 & 2\\ 1 &
        1\end{psmallmatrix}$.
  \item Deduce la estimación del error
        \[
        \Bigl|\frac{a_n}{b_n} - \sqrt2\Bigr|
        = \frac{1}{b_n\,(a_n + \sqrt2\,b_n)}
        \leq \frac{1}{2b_n^2},
        \]
        y prueba que decrece geométricamente con razón $3 -
        2\sqrt2$ \emph{(halla los \omterm{def:b2:reduction:eigen}{valores propios} de $M$ y el
        crecimiento de $b_n$)}.
  \item (Crecimiento general) A partir de la pregunta 9, demuestra:
        (a) si toda raíz cumple $\abs{\lambda_i} \leq \rho$,
        entonces $\abs{u_n} \leq C\,n^{m-1}\rho^n$ con $m = \max_i
        m_i$; (b) si hay una única raíz $\lambda_1$ de módulo
        máximo y $Q_1 \neq 0$, entonces $u_{n+1}/u_n \to
        \lambda_1$; compruébalo sobre la solución de la pregunta
        10.
\end{enumerate}

\medskip
\noindent\textbf{Parte IV --- Contar caminos y palabras.} Para un
grafo finito con conjunto de vértices $\{1, \dots, N\}$, la
\emph{matriz de adyacencia} $A$ tiene $A_{ij} = 1$ si $ij$ es una
arista y $0$ en caso contrario.
\begin{enumerate}[resume]
  \item Demuestra que $(A^n)_{ij}$ es el número de caminos de
        longitud $n$ de $i$ a $j$ (sucesiones de $n$ aristas, cada
        paso a lo largo de una arista).
  \item (El triángulo) Para el grafo completo de $3$ vértices, $A =
        J - I$: usando el \omterm{def:b2:reduction:eigen}{espectro} de $J$
        (\cref{ex:b2:linalg:onesmatrix}), prueba que
        \[
        (A^n)_{ii} = \frac{2^n + 2(-1)^n}{3},
        \qquad
        (A^n)_{ij} = \frac{2^n - (-1)^n}{3} \quad (i \neq j),
        \]
        y comprueba ambas fórmulas para $n = 2$ enumerando caminos.
  \item (Palabras sin $11$) Sea $w_n$ el número de palabras
        binarias de longitud $n$ sin dos $1$ consecutivos.
        Codifica las palabras por su última letra para obtener una
        matriz de transferencia, prueba que $w_{n+2} = w_{n+1} +
        w_n$, deduce $w_n = F_{n+2}$ (Fibonacci,
        \cref{exo:b2:reduction:5}) y da la tasa de crecimiento
        $\lim w_{n+1}/w_n$.
  \item (El camino) Para el grafo camino $1 - 2 - 3$, prueba que
        los \omterm{def:b2:reduction:eigen}{valores propios} de $A$ son $\sqrt2, 0, -\sqrt2$, con
        \omterm{def:b2:reduction:eigen}{vectores propios} $(1, \pm\sqrt2, 1)$ y $(1, 0, -1)$, y
        deduce que el número de caminos de longitud $n$ de un
        extremo al otro es $\bigl((\sqrt2)^n +
        (-\sqrt2)^n\bigr)/4$: cero para $n$ impar y $2^{\,n/2 - 1}$
        para $n$ par. Compruébalo para $n = 4$.
  \item (Fórmula de la traza) Prueba que el número total de caminos
        cerrados de longitud $n$ (con todos los puntos de partida)
        es $\operatorname{tr}(A^n) = \sum_i \lambda_i^n$, y
        verifícalo en el triángulo.
\end{enumerate}

\medskip
\noindent\textbf{Parte V --- Un anillo de sucesiones:
diagonalización simultánea.} Fijemos $k \geq 3$, sea $\omega =
\eu^{2\iu\pi/k}$ y sea $W \in \mathcal{M}_k(\C)$ el desplazamiento
\omterm{def:b2:structures:generated}{cíclico}: $W e_i = e_{i+1}$ (índices módulo $k$, columnas indexadas
$0, \dots, k-1$).
\begin{enumerate}[resume]
  \item Prueba que $W^{\mathsf T}$ es la matriz compañera de $X^k -
        1$, deduce que $\chi_W = \mu_W = X^k - 1$ y que $W$ es
        \omterm{def:b2:reduction:diag}{diagonalizable} con los $k$ \omterm{def:b2:reduction:eigen}{valores propios} simples
        $\omega^j$ y \omterm{def:b2:reduction:eigen}{vectores propios} $f_j = (1, \omega^{-j},
        \omega^{-2j}, \dots, \omega^{-(k-1)j})^{\mathsf T}$.
  \item Una matriz \emph{circulante} es $C = c_0 I + c_1 W + \dots
        + c_{k-1}W^{k-1}$. Prueba que todas las circulantes
        conmutan, que la base $(f_0, \dots, f_{k-1})$ las
        diagonaliza \emph{todas} simultáneamente y que los
        \omterm{def:b2:reduction:eigen}{valores propios} de $C$ son $\widehat c(\omega^j) = \sum_m
        c_m \omega^{jm}$, $j = 0, \dots, k-1$.
  \item Deduce que $\det C = \prod_{j=0}^{k-1} \widehat
        c(\omega^j)$, y comprueba que $k = 3$ recupera la
        factorización del~\cref{exo:b2:linalg:8}.
  \item (La media en el collar) Sea $x^{(n+1)} = Mx^{(n)}$ con $M =
        \frac12(W + W^{-1})$: cada uno de los $k$ números
        dispuestos en anillo se sustituye por la media de sus dos
        vecinos. Prueba que los \omterm{def:b2:reduction:eigen}{valores propios} de $M$ son
        $\cos(2\pi j/k)$ y que el coeficiente de $x^{(0)}$ sobre
        $f_0$ es la media $\frac1k\sum_m x^{(0)}_m$ \emph{(suma las
        coordenadas de los $f_j$)}.
  \item Concluye: para $k$ impar, $x^{(n)}$ converge al vector
        constante cuyo valor es la media de los valores iniciales;
        para $k = 4$, exhibe el \omterm{def:b2:reduction:eigen}{valor propio} responsable de la no
        convergencia y la obstrucción exacta (un coeficiente de
        media alternada que debe anularse).
  \item (Síntesis) En una frase cada uno: cómo la matriz compañera
        convierte el análisis de $(\mathcal R)$ en reducción; dónde
        el lema de descomposición en núcleos no necesitó dimensión
        finita; por qué los \omterm{def:b2:reduction:eigen}{valores propios} dominantes gobiernan las
        tasas de crecimiento y el error diofántico; por qué las
        potencias de la matriz de adyacencia cuentan caminos; y qué
        se gana con matrices que conmutan. Nombra las dos cumbres:
        el teorema fundamental de las recurrencias lineales y ---en
        el volumen del tercer año, para las matrices positivas de
        la parte IV--- el teorema de Perron--Frobenius.
\end{enumerate}
\end{problem}
